\begin{block}{Abstract}
    The computation of the topological susceptibility in 
    lattice gauge theory requires a smoothing procedure for 
    the gluonic topological charge, with cooling and gradient 
    flow being the most commonly adopted approaches. While 
    cooling is computationally cheaper, gradient flow admits 
    a Symanzik expansion that is easier to compute and allows 
    for a clean identification of leading order cutoff effects.
    We study these effects in pure SU(3) Yang--Mills theory using 
    ensembles generated with different gauge actions at lattice 
    spacings ranging from 0.04 fm to 0.12 fm. By combining different 
    levels of improvement for the flow kernel with several discretisations 
    of the topological charge operator, we study the contributions 
    to lattice artefacts arising from the action, flow, and observable,
    aiming to disentangle their respective effects. We present 
    preliminary results on the scaling of the topological susceptibility
    and the effectiveness of improvement strategies.
\end{block}